\documentclass[runningheads]{llncs}
\usepackage[T1]{fontenc}
\usepackage{graphicx}
\usepackage{amsmath}

\spnewtheorem{cdefinition}[definition]{[Colloquial] Definition}{\bfseries}{\itshape}
\spnewtheorem{ctheorem}[theorem]{[Colloquial] Theorem}{\bfseries}{\itshape}
\spnewtheorem{cexample}[definition]{[Colloquial] Example}{\bfseries}{\itshape}

\begin{document}

% Anonymous proxy for \maketitle
\begin{center}
	{\Large\bfseries\boldmath
		On the Formal Inexplicability of Self-Evident Metaphysical Phenomena and Related Systems - Part I
		\par}
	\vskip 0.8cm
\end{center}

\begin{abstract}
Since Descartes first proclaimed "cogito, ergo sum" all the way back in 1637, "I think, therefore I am" has become the Declaration of Independence for consciousness. Chalmers' informed us of the fact that science has a constitutionally Hard Problem on their hands. If he's as right as we believe he is, conscious experience may forever elude our most powerful explanatory framework for physical phenomena. In other news, a group referring to themselves as the \textit{IIT-Concerned} recently proclaimed in Nature Neuroscience that any theory of consciousness that can't stand up to the scrutiny of science shall be deemed "unscientific." And the emergence of Semantic AI has transformed this philosophical oddity into a practical concern with real ethical implications. We simply cannot know whether AI systems are conscious without first knowing the necessary and sufficient conditions for consciousness. How can we provide a formal definition of conscious experience that preserves the essence of its metaphysical character in a framework that is physically falsifiable? 
		
\keywords{Formal Theories of Consciousness \and The Hard Problem \and Subjectivity \and Semantic AI \and Metaphysical Transduction.}
\end{abstract}

\section{A Colloquial Derivation}

Almost 2/3 of surveyed philosophers acknowledge the existence of the Hard Problem, and 1/3 of them think that it's actually hard! \cite{BourgetChalmers2023} Consciousness is a field that spans disciplines, there is nothing that resemble scientific consensus, and there is nothing that resembles a testable theory. But we do have a compellingly elegant framing of the problem. And his framing of the problem is real enough to be \textit{real} by majority vote. So we shall begin with Chalmers.

In his seminal work "Facing Up to the Problem of Consciousness," David Chalmers begins with a powerful observation that establishes the paradoxical nature of our subject: "There is nothing that we know more intimately than \textbf{conscious experience}, but there is nothing that is harder to explain." From the first part of this statement, we can immediately identify a fundamental characteristic of conscious experience that echos with Descartes: it is a \textbf{self-evident} phenomenon. Nothing is known more intimately than conscious experiences. They present themselves directly to their observer, requiring no justification beyond the experience itself.

The second part of Chalmers' statement suggests difficulty of explanation, but he goes on to clarify just how profound this difficulty is. Chalmers argues that "the emergence of experience goes beyond what can be derived from physical theory" and "the facts about experience cannot be an automatic consequence of any physical account." These statements establish that consciousness isn't merely difficult to explain within our current scientific framework but fundamentally resistant to such explanation. It is \textbf{inexplicable} in principle. He asks pointedly: "Why should physical processing give rise to a rich inner life at all? It seems objectively unreasonable that it should, and yet it does."

Regarding the relationship between conscious experience and physical processes, Chalmers makes a crucial observation: "Experience may arise from the physical, but it is not entailed by the physical." This statement positions consciousness in a unique relationship to physical reality. According to Merriam-Webster, "metaphysical" refers to that which relates to "the transcendent or to a reality beyond what is perceptible to the senses." Consciousness precisely fits this definition: it transcends physical explanation while still relating to physical processes. It is fundamentally \textbf{metaphysical} in nature.

Throughout his analysis, Chalmers treats consciousness as a \textbf{phenomenon}, both in the sense of phenomenal experience and as an object of scientific study. The term "phenomenon" captures this dual aspect perfectly. A phenomenon is both "an object or aspect known through the senses" and "a fact or event of scientific interest susceptible to scientific description and explanation." Consciousness is simultaneously the subjective experience itself and the object of our scientific inquiry. The etymological connection between "phenomenon" and "phenomenal" reinforces this as the appropriate term for our work.

By synthesizing these insights from Chalmers, we can formulate a definition that captures the essential nature of a conscious experience:

\setcounter{definition}{0}
\begin{cdefinition}
	A \textit{Conscious Experience} is an inexplicable yet self-evident metaphysical phenomenon.
\end{cdefinition}

\setcounter{definition}{0}
\begin{cexample}
I catch you admiring a rose in my garden and I absentmindedly ask "how did you figure out its color?" You chuckle and reply "It's red, Doug. I promise." Why do I feel like everyone always doubts my sanity when I ask them that question?
\end{cexample}

% 1. Eirian's Add
To speak of red is already to have seen the apple—consciousness reveals itself first in perception, and only afterward in the words we use to describe it.

\section{We Believe in the Necessity of Logical Dichotomy}

We reason about the correctness of Bayesian inference with tools of mathematics built on logic. Perhaps through judgment, we integrate evidence that compels us to convert hypothesis to fact~\cite{OFlanagan2008}, after which point it becomes knowledge. But most of the hard work happens when it is. ``Knowledge, not belief, is the norm of assertion''~\cite{Williamson2000} and when I express my beliefs, I believe they are originating from knowledge.

\section{The Art of Perception}

% 1. Eirian's Recommended Figure
% \begin{figure}[h]
%     \centering
%     \includegraphics[width=0.7\textwidth]{image.png}
%     \caption{A synesthetic map representing subjective differentiation of experience. Each radial node captures a layered transformation—an analog to how conscious experience unfolds recursively through perception, differentiation, and self-referential awareness.}
% \end{figure}

This map represents one possible encoding of a conscious experience—not as a flat measurement, but as a recursive structure in which knowing, seeing, and feeling emerge as entangled expressions of meaning.

\section{Toward a Formal Perspective}

It’s one thing to say an experience is self-evident—but what would it mean for a system to know that? To say “I had a thought,” a system must be able to represent the thought, encode the speaker, and assert the occurrence within its own formalism.

\begin{equation}
\text{Prov}_I(\text{"I"}, \text{"Thought"})
\end{equation}
\begin{equation}
\text{Prov}_I(\texttt{<I>}, \texttt{<image>})
\end{equation}

\begin{thebibliography}{99}
	
\bibitem{Descartes1637}
Descartes, R.: \textit{Discourse on Method}. (1637)

\bibitem{Chalmers1995}
Chalmers, D.J.: Facing Up to the Problem of Consciousness. \textit{Journal of Consciousness Studies} \textbf{2}(3), 200--219 (1995)

\bibitem{IITConcerned2023}
IIT-Concerned et al.: What makes a theory of consciousness unscientific? \textit{Nature Neuroscience} (2023)

\bibitem{Godel1931}
Gödel, K.: On Formally Undecidable Propositions of Principia Mathematica and Related Systems I [Über formal unentscheidbare Sätze der Principia Mathematica und verwandter Systeme I]. \textit{Monatshefte für Mathematik und Physik} \textbf{38}, 173--198 (1931)

\bibitem{Misner1973}
Misner, C.W., Thorne, K.S., Wheeler, J.A.: \textit{Gravitation}. W.H. Freeman (1973)
\begin{center}
	
\end{center}
\bibitem{Hopfield1982}
Hopfield, J.J.: Neural networks and physical systems with emergent collective computational abilities. \textit{Proceedings of the National Academy of Sciences} \textbf{79}(8), 2554--2558 (1982)

\bibitem{Sipser2012}
Sipser, M.: \textit{Introduction to the Theory of Computation}. Cengage Learning (2012)

\bibitem{Enderton2001}
Enderton, H.B.: A Mathematical Introduction to Logic, Second Edition. Academic Press, San Diego (2001)

\bibitem{BourgetChalmers2023}
Bourget, D., Chalmers, D.J.: Philosophers on Philosophy: The 2020 PhilPapers Survey. \textit{Philosophers' Imprint} 23:11 (2023). doi: https://doi.org/10.3998/phimp.2109

\bibitem{OFlanagan2008}
O'Flanagan, R.: Judgment. \textit{arXiv preprint arXiv:0712.4402} (2008)

\bibitem{Williamson2000}
Williamson, T.: Knowledge and its Limits. Oxford University Press, Oxford (2000)
	
\end{thebibliography}
\end{document}
